\subsection{APPLICATION TO THE SUBSPACE ASSOCIATED WITH A p-VECTOR}

Let K be a field and E a vector space over K. Recall that with every p-vector \(z \in \bigwedge^{p}(\mathbf{E})\) is associated a finite-dimensional subspace \(M_{z}\) of E, namely the smallest vector subspace M of E such that \(z \in \bigwedge^{p}(\mathbf{M})\) (§ 7, no. 2, Corollary to Proposition 4).

PROPOSITION 14. (i) The orthogonal of \(\mathbf{M}_z\) in \(\mathbf{E}^*\) is the set of \(x^{*}\in \mathbf{E}^{*}\) such that \(x^{*}\lrcorner z = 0\) .

\begin{enumerate}
\def\labelenumi{(\roman{enumi})}
\setcounter{enumi}{1}
\tightlist
\item
  The subspace \(\mathbf{M}_z\) associated with \(z\) is the image of \(\bigwedge^{p-1}(\mathbf{E}^*)\) under the mapping \(\lambda_z: u^* \mapsto u^* \lrcorner z\) of \(\bigwedge^{p-1}(\mathbf{E}^*)\) into \(\mathbf{E}\) .
\end{enumerate}

Let \(\mathbf{N}\) denote the image of \(\lambda_z\) . For \(x^* \in \mathbf{E}^*\) and \(u^* \in \bigwedge^{p-1}(\mathbf{E}^*)\) ,

\[
\begin{array}{r l} \langle \theta_ {\wedge} (x ^ {*}), u ^ {*} \lrcorner z \rangle = \langle \theta_ {\wedge} (u ^ {*} \wedge x ^ {*}), z \rangle & = (- 1) ^ {p - 1} \langle \theta_ {\wedge} (x ^ {*} \wedge u ^ {*}), z \rangle \\ & = (- 1) ^ {p - 1} \langle \theta_ {\wedge} (u ^ {*}), x ^ {*} \lrcorner z \rangle . \end{array}
\]

Therefore, for \(x^*\) to be orthogonal to \(\mathbf{N}\) , it is necessary and sufficient that \(x^* \lrcorner z\) be orthogonal to \(\theta_\wedge (\bigwedge (\mathrm{E}^*))\) . Now, the latter condition is equivalent to saying that \(x^* \lrcorner z = 0\) ; for let \((e_\lambda)_{\lambda \in L}\) be a basis of \(\mathbf{E}\) ; giving \(L\) a total ordering, it has been seen (§ 7, no. 8, Theorem 1) that the \(e_j\) , for \(J\) running through the set \(\mathfrak{F}(L)\) of finite subsets of \(L\) , form a basis of \(\bigwedge (\mathrm{E})\) ; it then follows from formula (30) of no. 5 that the elements \(\theta_{\wedge}(e_{\mathrm{J}}^{*})\) are, to within a sign, the coordinate forms on \(\bigwedge(\mathrm{E})\) relative to the basis \((e_{\mathrm{J}})\) ; whence our assertion.

The orthogonal of N therefore consists of the \(x^{*} \in E^{*}\) such that \(x^{*} \lrcorner z = 0\) and the conclusion of (i) will therefore follow from (ii).

We show first that \(\mathbf{N} \subset \mathbf{M}_z\) . Let \(\mathbf{M}\) be a vector subspace of \(\mathbf{E}\) such that \(z \in \bigwedge (\mathbf{M})\) and let \(j: \mathbf{M} \to \mathbf{E}\) be the canonical injection; let \(\mu_z\) denote the mapping \(v^* \mapsto v^* \lrcorner z\) of \(\bigwedge^{p-1}(\mathbf{M}^*)\) into \(\mathbf{M}\) ; it follows from formula (60) of no. 7 that there is a canonical factorization

\[
\lambda_ {z}: \bigwedge^ {p - 1} (\mathrm{E} ^ {*}) \xrightarrow {\bigwedge^ {p - 1} ({} ^ {t} j)} \bigwedge^ {p - 1} (\mathrm{M} ^ {*}) \xrightarrow {\mu_ {z}} \mathrm{M} \xrightarrow {j} \mathrm{E}
\]

which proves that \(\mathbf{N} \subset \mathbf{M}\) and hence \(\mathbf{N} \subset \mathbf{M}_z\) by definition of \(\mathbf{M}_z\) . It remains to verify that \(\mathbf{N} = \mathbf{M}_z\) . Suppose the converse: there would then exist a basis \((e_i)_{1 \leqslant i \leqslant n}\) of \(\mathbf{M}_z\) and an element \(x^* \in \mathbf{E}^*\) such that \(\langle x^*, e_1 \rangle = 1\) , \(\langle x^*, e_j \rangle = 0\) for \(2 \leqslant j \leqslant n\) and such that \(x^*\) is orthogonal to \(\mathbf{N}\) and hence \(x^* \lrcorner z = 0\) . We write \(z = \sum_{\mathbf{H}} a_{\mathbf{H}} e_{\mathbf{H}}\) , where the sum is taken over the subsets of \([1, n]\) with \(p\) elements. By (68) (no. 9),

\[
\begin{array}{l l} x ^ {*} \lrcorner e _ {\mathrm{H}} = 0 & \text { if } \quad 1 \notin \mathrm{H} \\ x ^ {*} \lrcorner e _ {\{1 \} \cup \mathrm{H}} = e _ {\mathrm{H}} & \text { if } \quad \mathrm{H} \subset [ 2, n ] \end{array}
\]

which shows that the relation \(x^{*} \lrcorner z = 0\) implies \(a_{\mathbf{H}} = 0\) for \(1 \in \mathbf{H}\) . But this is impossible, for \(z\) would then belong to \(\bigwedge^{p}(\mathbf{M}')\) , where \(\mathbf{M}'\) is the subspace of \(\mathbf{M}\) generated by \(e_2, \ldots, e_n\) .

